\chapter{Objetivos}
\label{cap:objetivos}

\section{Objetivo geral}

Desenvolver e documentar o Cardly, aplicativo móvel de flashcards com revisão espaçada, demonstrando integração completa entre React Native (Expo/TypeScript) e Spring Boot (Java), incluindo autenticação, CRUD, comunicação REST, organização em camadas, experiência de usuário responsiva e preparação para ambiente de produção.

\section{Objetivos específicos}

\begin{enumerate}
    \item Implementar fluxos de autenticação com credenciais locais e Google SSO, utilizando JWT stateless \cite{jwt_io,google_oauth,spring_security_jwt}.
    \item Modelar assuntos (\textit{decks}), cartas e respostas com regras de dificuldade e intervalos entre tentativas definidos pelo sistema.
    \item Disponibilizar dashboard com indicadores de progresso e calendário de dias com atividade de estudo, inspirado em aplicativos de acompanhamento diário \cite{clue_app}.
    \item Expor API REST paginada com endpoints \texttt{/search} baseados em Specification, sem \texttt{getAll} genérico \cite{jpa_specification}.
    \item Construir front-end com NativeWind, modais de confirmação e toasts equilibrados \cite{nativewind_docs,toast_message_rn}.
    \item Oferecer painel exclusivo de superadministrador para gestão de usuários, assuntos e cartas.
    \item Validar o sistema com testes automatizados e evidências visuais para arguição técnica \cite{junit5_spring,jest_testing_library}.
    \item Preparar pipeline de deploy (Railway/Vercel ou equivalente) para demonstração funcional \cite{railway_docs,vercel_docs,expo_eas_build}.
\end{enumerate}

\section{Mapeamento aos critérios de avaliação do professor}

A Tabela~\ref{tab:criterios-avaliacao} relaciona explicitamente os critérios divulgados na plataforma da disciplina (conforme \texttt{todo.md} do projeto) às seções deste relatório e às evidências implementadas no Cardly. Esse mapeamento orienta a apresentação prática e a arguição técnica.

\begin{table}[htbp]
    \caption{Mapeamento entre critérios de avaliação, seções do relatório e evidências no Cardly}
    \label{tab:criterios-avaliacao}
    \centering
    \small
    \begin{tabular}{p{4.2cm} p{3.8cm} p{6.5cm}}
        \toprule
        \textbf{Critério (disciplina)} & \textbf{Seção do relatório} & \textbf{Evidência no projeto} \\
        \midrule
        Integração Front-End (RN/Expo/TS) e Back-End (Spring Boot) &
        Cap.~\ref{cap:arquitetura}, \ref{cap:backend}, \ref{cap:frontend} &
        Módulos \texttt{*Api.ts} no app; controllers REST no backend; JSON paginado \\
        \addlinespace
        Arquitetura cliente-servidor &
        Cap.~\ref{cap:arquitetura} &
        Diagramas de camadas; API stateless; PostgreSQL centralizado \\
        \addlinespace
        Autenticação, CRUD e comunicação com API &
        Cap.~\ref{cap:seguranca}, \ref{cap:backend}, \ref{cap:frontend} &
        JWT + Google SSO; endpoints protegidos; CRUD de decks/cartas/usuários \\
        \addlinespace
        Organização de código e separação de responsabilidades &
        Cap.~\ref{cap:backend}, \ref{cap:frontend} &
        Padrão 1.0/1.1 no backend; screens, components, api e auth no front \\
        \addlinespace
        UI/UX, responsividade e experiência &
        Cap.~\ref{cap:frontend} &
        FlipCard animado; paleta 2.0; toasts e modais; dashboard \\
        \addlinespace
        Preparo para produção &
        Cap.~\ref{cap:deploy} &
        Variáveis de ambiente; build EAS; deploy Railway; CORS \\
        \addlinespace
        Funcionamento completo da aplicação &
        Cap.~\ref{cap:testes}, \ref{cap:conclusao} &
        Testes de serviço; roteiro de demonstração; prints \\
        \addlinespace
        Segurança adequada &
        Cap.~\ref{cap:seguranca} &
        BCrypt, JWT, roles, SecureStore, OWASP mobile \\
        \addlinespace
        Clareza na apresentação e domínio técnico &
        Cap.~\ref{cap:objetivos}, \ref{cap:conclusao} &
        Mapeamento de critérios; glossário implícito; decisões justificadas \\
        \addlinespace
        Estabilidade geral &
        Cap.~\ref{cap:testes} &
        Testes unitários backend; tratamento de erros; validações \\
        \bottomrule
    \end{tabular}
    \\[0.5em]
    \footnotesize Fonte: elaborado pelos autores com base no enunciado da disciplina e em \texttt{todo.md}.
\end{table}

\section{Mapeamento à metodologia e atividades da apresentação}

A Tabela~\ref{tab:atividades-apresentacao} indica como cada atividade exigida na metodologia será conduzida na banca.

\begin{table}[htbp]
    \caption{Atividades de apresentação e correspondência no Cardly}
    \label{tab:atividades-apresentacao}
    \centering
    \small
    \begin{tabular}{p{5.5cm} p{8.5cm}}
        \toprule
        \textbf{Atividade (professor)} & \textbf{Como será demonstrada} \\
        \midrule
        Proposta (problema e solução) & Slides iniciais + Cap.~1 deste relatório \\
        Fluxo completo (login, navegação, CRUD, API) & Demo no emulador/dispositivo: Login $\rightarrow$ Decks $\rightarrow$ Estudo $\rightarrow$ Dashboard \\
        API em funcionamento & Postman ou log ao vivo; endpoints \texttt{/api/auth}, \texttt{/api/decks}, \texttt{/api/cards} \\
        Arquitetura (Controller/Service/Repository/Hooks/Services) & Diagramas Cap.~\ref{cap:arquitetura}; navegação de pastas nos repositórios \\
        UI/UX e layout & FlipCard, paleta de cores, calendário no dashboard \\
        Build ou preparação para deploy & \texttt{eas build} ou APK pré-gerado; URL Railway ativa \\
        Questionamentos técnicos & Equipe dividida por domínio (ver roteiro de apresentação em material complementar) \\
        \bottomrule
    \end{tabular}
\end{table}

\section{Delimitações}

Não fazem parte do escopo mínimo desta entrega final de documentação: implementação completa do módulo de amizades (solicitações entre usuários), embora conste nos requisitos funcionais planejados; geração automática de slides \texttt{pptx}; e revisão por três agentes revisores exigida internamente pela equipe. Esses itens constam como pendências na conclusão deste relatório.
